\documentclass[11pt,a4paper]{article}
\usepackage[utf8]{inputenc}
\usepackage[T1]{fontenc}
\usepackage[backend=biber, style=numeric]{biblatex}
\addbibresource{refs.bib}

\usepackage{centrale}

\hypersetup{
    pdftitle={Titre à insérer},
    pdfauthor={Auteur à insérer},
    pdfsubject={Sujet à insérer},
    pdfproducer={Conversion PDF à insérer},
    pdfkeywords={Quelques mots-clés à insérer} %
}

\DeclareGraphicsRule{.ai}{pdf}{.ai}{} % pour insérer des documents .ai
\graphicspath{ {./img/} {./eps/}} % pour ne pas avoir à ajouter eps/ton-image.jpg

% ------------- Packages spéciaux, nécessaires pour ce rapport, à insérer ici ------------- 


\begin{document}

% --------------------------------------------------------------
%                       Page de garde
% --------------------------------------------------------------

\begin{titlepage}
\begin{center}


{\large 4\up{ème} CEI d'Informatique Quantique avec CA-CIB}\\[0.5cm]

{\large Dominante INFONUM - Projet de fin d'études}\\[0.5cm]

% Title
\rule{\linewidth}{0.5mm} \\[0.4cm]
{ \huge \bfseries Quantum Embedding for Advanced Analytics in Finance \\[0.4cm] }
\rule{\linewidth}{0.5mm} \\[0.4cm]


\begin{figure}[h]
    \centering
    \begin{minipage}{0.6\textwidth}
        \centering
        \includegraphics[width=\textwidth]{./logos/cs_logo.jpg}
    \end{minipage}%
    \hfill
    \begin{minipage}{0.4\textwidth}
        \centering
        \includegraphics[width=\textwidth]{./logos/cacib_logo.png}
    \end{minipage}
\end{figure}

% Author and supervisor
\noindent
\begin{minipage}{0.4\textwidth}
  \begin{flushleft} \large
    \emph{Auteurs :}\\
    M\up{me} Malek \textsc{Bouhadida}\\
    M. Pierre \textsc{El Anati}\\
    M. Pierre \textsc{Jourdin}\\
  \end{flushleft}
\end{minipage}%
\begin{minipage}{0.6\textwidth}
  \begin{flushright} \large
    \emph{Encadrants CentraleSupélec:} \\
    M.~Damien \textsc{Rontani}\\
    M.~Stéphane \textsc{Vialle}\\
    \vspace{0.5cm}
    \emph{Encadrants Crédit Agricole CIB:} \\
    M.~Frédéric \textsc{Bernard}\\
    M.~David \textsc{Dubuis}\\
    M.~Christophe \textsc{Michel}\\
    M.~Achraf \textsc{Seddik}
  \end{flushright}
\end{minipage}

\vfill

\begin{center}
    {\large Année Universitaire 2024-2025}
\end{center}

\end{center}
\end{titlepage}

% --------------------------------------------------------------
%                            Abstract
% --------------------------------------------------------------
\newpage
\thispagestyle{empty}

\vspace*{\fill}
\noindent\rule[2pt]{\textwidth}{0.5pt}\\
{\textbf{Résumé :}}
\lipsum[1]

{\noindent\textbf{Mots clés :}}
Lorem ipsum dolor sit amet, consectetur adipiscing elit. Sed non risus. Suspendisse lectus tortor.
\\
\noindent\rule[2pt]{\textwidth}{0.5pt}
\begin{center}
  École centrale de Lyon\\
  36, Avenue Guy de Collongue\\
  code du labo concerné\\
  69134 Écully
\end{center}
\vspace*{\fill}

% --------------------------------------------------------------
%                    Table des matières 
% --------------------------------------------------------------
\newpage
\thispagestyle{empty}
\tableofcontents
%\newpage

% --------------------------------------------------------------
%                         Début du corps
% --------------------------------------------------------------
\newpage
\section{Introduction}
Le projet est réalisé en collaboration avec CA-CIB (Crédit Agricole Corporate and Investment
Bank) et le CEI CentraleSupélec. Dans un contexte de veille technologique, CA-CIB souhaite explorer le potentiel de l’informatique quantique pour le traitement des
données financières.
L’objectif principal est de développer des algorithmes
d’embedding exploitant l’informatique quantique. Ces algorithmes sont destinés à fournir des
représentations efficaces des données financières, dans le but de préparer la banque à l’éventuelle
démocratisation des ordinateurs quantiques. En effet, si ces technologies deviennent opérationnelles
dans le futur, CA-CIB souhaite disposer d’une base d’algorithmes prête à être déployée
pour en tirer parti immédiatement.
Le problème posé peut être formulé ainsi : Comment utiliser l’informatique quantique pour
générer des embeddings optimisés sur des données financières en anticipant les évolutions technologiques
à venir ?



% --------------------------------------------------------------
%                         Etat de l'art
% --------------------------------------------------------------

\newpage
\section{État de l'art}
Pour établir un état de l’art solide, nous avons commencé par une sélection d’articles en nous
basant sur leurs abstracts. Cette étape préliminaire nous a permis d’identifier les travaux les plus
cités et alignés avec notre projet, sur lesquels nous avons approfondi nos recherches.
\subsection{Apprentissage automatique classique}
\subsection{Apprentissage automatique quantique}




\lipsum[1] 

Comme dans la figure 
\reff{fig:id-de-la-figure}.


\begin{figure}[ht!]
    \centering
    \includegraphics[width=0.3\textwidth]{example-image-a}
    \caption{Insérer ici le sous-titre.}
    \label{fig:id-de-la-figure}
\end{figure}

%\begin{listing}[ht]
%\inputminted[mathescape,
               %numbersep=5pt,
               %gobble=2,
               %frame=lines,
               %framesep=2mm]{python}{code/}
%\caption{Du code informatique}
%\label{listing:id-du-code}
%\end{listing}


% --------------------------------------------------------------
%                            Schéma d'ensemble
% --------------------------------------------------------------
\newpage
\section{Les différentes approches explorées}
\subsection{Implémentation de l'article \og Quantum embeddings for machine learning \fg}
\subsection{QLSTM}
Pour faire suite à l'état de l'art, nous avons décidé de nous orienter vers l'article "Quantum Long Short-Term Memory" [Yen-Chi Chen \emph{et al}., 2020], qui propose un modèle LSTM hybride quantique-classique appelé QLSTM. L'article montre que le modèle proposé apprend avec succès plusieurs types de données temporelles, où pour certains cas de test, la version QLSTM converge plus rapidement, ou, de manière équivalente, atteint une meilleure précision, que son homologue classique.\\L'architecture est donc celle d'un LSTM classique comme preséntée dans la figure 1, à la différence où les couches de neurones cachées sont remplacées par des \emph{Variational Quantum Circuits (VQC)}. Son fonctionnement est le suivant : les entrées du QLSTM sont des données tabulaires au format classique. Ces entrées vont ensuite être passées dans un VQC qui va faire l'embedding des données classiques en qubits afin que celles-ci puissent être computable dans le circuit quantique.\\Cela se réalise en passant chacune des features dans une porte de Haddamard puis en faisant une rotation d'angle \( \theta_{1} \) autour de l’axe des ordonnées, suivie d’une rotation d’angle \( \theta_{2} \) autour de l’axe des abscisses. Cette suite d’opérations est appelée \textit{Amplitude Encoding}.
Une fois les \textit{features} encodées sous forme de qubits, ces derniers sont injectés dans un circuit quantique composé d’opérations d’intrication et de superposition, permettant d’exploiter les propriétés fondamentales de la mécanique quantique pour le traitement de l’information. Enfin chaque qubit est mesuré afin d'obtenir un scalaire qui sera ensuite passé à la fonction d'activation du circuit classique.
\subsection{QuLTSF}
\subsection{Une autre sous-partie}
\lipsum[2]

% --------------------------------------------------------------
%                            Organisation
% --------------------------------------------------------------
\newpage
\section{Organisation}
\subsection{Une autre sous-partie}
\lipsum[2]

% --------------------------------------------------------------
%                            Résultats obtenus
% --------------------------------------------------------------
\newpage
\section{Résultats}
\subsection{Une autre sous-partie}
\lipsum[2]

% --------------------------------------------------------------
%                            Conclusion
% --------------------------------------------------------------
\newpage
\section{Conclusion}


\lipsum[1]

\section{Test All References}
Testing all citations in order:

\begin{enumerate}
    \item Dynamic Network Embedding: \cite{dynamic_network_embedding_survey}
    \item ASNE Algorithm: \cite{asne_social_network}
    \item STOCK2VEC: \cite{stock2vec}
    \item Word2Vec: \cite{word2vec}
    \item Word Embedding Survey: \cite{word_embedding_survey}
    \item Stochastic Neighbor Embedding: \cite{stochastic_neighbor_embedding}
    \item GloVe: \cite{glove}
    \item Time2Vec: \cite{time2vec}
    \item Stock Embeddings: \cite{stock_embeddings}
    \item Attention Is All You Need: \cite{attention_is_all_you_need}
    \item Price Graphs: \cite{price_graphs}
    \item T-REP: \cite{trep}
    \item Adaptive Embedding: \cite{adaptive_embedding}
    \item Time Series 64 Words: \cite{timeseries_64words}
    \item Quantum ML Finance: \cite{quantum_ml_finance}
    \item Quantum Embeddings ML: \cite{quantum_embeddings_ml}
    \item Neural Quantum Embedding: \cite{neural_quantum_embedding}
    \item Quantum Embedding Search: \cite{quantum_embedding_search}
    \item Ansatz Design Review: \cite{ansatz_design_review}
    \item Data-driven Quantum Embedding: \cite{data_driven_quantum_embedding}
    \item QuLTSF: \cite{qultsf}
    \item Quantum LSTM: \cite{quantum_lstm}
    \item ARIMAX Thailand: \cite{arimax_thailand}
    \item Hybrid ARIMA-NN: \cite{hybrid_arima_nn}
\end{enumerate}

\printbibliography[title=Complete Bibliography]

\end{document}